\section{总结}
本文围绕赛题二的 3D 异构 CIM 资源调度问题，完成了从模型与轨迹解析、Weight-Cube 切分、共现感知映射到静态调度和独立验证的可运行链路。V5 primary 在严格容量约束下将平衡 mock trace 的模拟延迟降低 \PrimaryImprovement{}，并保持映射率 1.0；大规模 metadata-only 实验进一步说明链路可扩展到 \LargeExperts{} 专家。结果表明，将专家激活统计用于分组、副本与调度决策，能够在当前模拟设置中降低切换和流水线等待。

该结论的适用范围限定于项目的资源模型、mock/synthetic 轨迹和已记录的参数。下一步将以官方轨迹进行 parser smoke、主方案与 V5.1 候选的统一 replay，并结合可获得的 CIM 芯片文献或测量数据校准切换、z 层访问和跨 Sub-Cube 传输参数。
